\subsubsection*{Revision: communication accounting and adaptive baselines}
\textcolor{blue}{In response to the reviewers' concerns about constrained networks and fair empirical comparisons, we explicitly account for communication in all new review experiments.}
One gossip round consists of one local exchange over every directed neighbor relation induced by the graph.
If $d$ is the model dimension and $|E_{\mathrm{dir}}|$ is the number of directed non-self neighbor transmissions, then one gossip round communicates $|E_{\mathrm{dir}}|d$ real scalars for \algname{}-\pmb{ii}, DOGD, and the adaptive decentralized baselines, and $|E_{\mathrm{dir}}|(d+1)$ real scalars for \algname{}-\pmb{i} because it additionally communicates the scalar wealth.
Thus, with a constant schedule $q(t)=1$, the communication cost is $\mathcal{O}(T |E_{\mathrm{dir}}| d)$ for \algname{}-\pmb{ii}, whereas a linear schedule has $\mathcal{O}(T^2 |E_{\mathrm{dir}}| d)$ cost.
This is why our practical experiments emphasize the constrained regime $q(t)=1$ and report the theoretical linear schedule as a sufficient, not mandatory, condition for the current proof.

\textcolor{blue}{We also added decentralized adaptive-gradient baselines using the same gossip interface as DOGD: AdaGrad, RMSProp, Adam, AdamW, momentum, and Nesterov-type momentum.}
These methods are not parameter-free in the sense of our theoretical claims because they still require stepsize, momentum, and/or decay choices, but they are important controls for the reviewers' question of whether DECO's behavior is merely an effect of adapting the update magnitude to past gradients.
The added code reports cumulative network loss, average network loss, cumulative local loss, and total communicated scalars for each method.
